\documentclass[a4paper]{article}

\usepackage{graphicx}
\usepackage{amsmath,amssymb}
\usepackage{minted}
\usepackage{multirow}
\usepackage{float}
\usepackage{todonotes}
\usepackage{forest}
\usepackage{xstring}
\usepackage{xcolor}
\usepackage[most]{tcolorbox}
\usepackage{xparse}
\usepackage{natbib}

\usetikzlibrary{graphs,graphdrawing}
\usegdlibrary{layered}

% for external graphviz
\input{/home/donkarlo/Dropbox/repo/nd_ai_project/src/nd_ai/knoweldge_representation/language/formal/computer/programming/tex/latex/code_clips/external_graphviz.tex}

% for tags
\input{/home/donkarlo/Dropbox/repo/nd_ai_project/src/nd_ai/knoweldge_representation/language/formal/computer/programming/tex/latex/parts/control_sequence/kinds/semtag/semtag.tex}

% for links
\input{/home/donkarlo/Dropbox/repo/nd_ai_project/src/nd_ai/knoweldge_representation/language/formal/computer/programming/tex/latex/code_clips/seehere.tex}

\title{Autopoiesis}
\date{}

\begin{document}
    \maketitle
    
    Autopoiesis one of several current theories of life, refers to a system capable of producing and maintaining itself by creating its own parts
    \seeherefooter{https://en.wikipedia.org/wiki/Autopoiesis}
    
    \cite{maturana-varela-1980-autopoiesis-and-cognition-the-realization-of-the-living}
    defines for the first time a living being as a system that continuously produces and reproduces itself. Its parts generate the very network that, in turn, produces those same parts.
    So, being alive means having a self-producing and self-maintaining organization, not merely possessing a certain material or behavior.
    
    
    
    \bibliographystyle{plainnat}
    \bibliography{/home/donkarlo/Dropbox/repo/nd_shared_research/bibliography/bibliography.bib}
\end{document}